% ===== §11 =====
\section{The Wager, not the Rehearsal}
\label{p27:sec:wager}

\noindent
The parts to this point have set out what relational understanding is, what it does, and why, done legitimately, it is good. A reader might expect the paper now to describe how it is practiced, how one learns and rehearses the capacities it requires until they are secure, in the way one rehearses any skill. This part argues that intimacy resists such a description at its root, and that the resistance is a structural feature and not a practical inconvenience that follows from the framework's own conception of value and time. The capacities of relational understanding, in the domain of intimacy, are wagered and not rehearsed.

A rehearsal is possible where an action can be tried, its result observed, and the action adjusted and tried again, the feedback returning quickly enough and the trials repeating freely enough that the capacity is shaped by the loop of attempt and correction. Much can be learned in this way, and much of what surrounds relational understanding can be. The recognition of another's world, the entry into it, the labor of translation, admit of practice in settings where the loop is available, in ordinary conversation, in friendship, in the encounter with a difficult text or an unfamiliar discipline, in the professional formations of teaching and care. There the crossing of worlds can be attempted, its result seen with reasonable speed, and the attempt refined. These capacities are, in a sense the next section will make exact, near-practicable.

Intimacy is different, and the difference follows from what the framework holds about relational value and its time. Value, on this account, is generated in the crossing of worlds, and its generation is historical, so that what is generated at one moment may disclose its value only much later, when the field of relations has grown enough to observe it. The companion work drew this out for the observation of concepts, where a rendering's significance may become observable only once a later rendering arrives to throw it into relief, and where the field of relations through which a matter is observed is itself generated over time. In an intimate bond this temporal structure is at its most extended. A word given, a room made for the other's difficulty, a generative thing offered at a hard moment, may disclose its value only after years, when the bond has developed to the point from which what was given can at last be observed for what it was. The feedback on an act of relational understanding, in intimacy, may return only after a passage of time so long that it cannot serve as the correction in a loop of practice.

To this delay is joined a second feature, that the trials do not repeat under the same conditions and often do not repeat at all. An intimate bond is a single history, not a series of independent attempts. A crossing offered at a given moment is offered once, into a bond in the state that its whole prior history has produced, and it cannot be withdrawn and offered again to the same bond in the same state, for the offering has already altered the state. Where a rehearsal requires that the trial be repeatable, the intimate bond is constitutively unrepeatable, a single unfolding in which each act enters once and changes the ground of the next. The partner one addresses today is not the partner one will address tomorrow, for the bond will have moved, carried in part by today's address.

The delay of the feedback and the unrepeatability of the trial together remove the loop of attempt and correction on which rehearsal depends. One cannot, in the deepest reaches of an intimate bond, try a way of understanding the other, observe promptly whether it succeeded, and adjust and try again under the same conditions, because the observation may be long delayed and the conditions do not return. What one does instead, in offering an act of relational understanding into an intimate bond, is to give something whose value one cannot presently confirm, into a history that will not run again, on the trust that it is the right thing to give. This is the structure of a wager and not of a rehearsal, and it is the structure that sets the cultivation of these capacities in intimacy apart from their cultivation anywhere the loop is available.

\begin{claim}[The wager of intimacy]
In the domain of intimacy the capacities of relational understanding are wagered, not rehearsed. Because relational value is generated historically, the worth of an act of understanding may return only after a long passage of time, and because the bond is a single unrepeatable history, the act cannot be tried again under the same conditions. The loop of attempt and prompt correction on which rehearsal depends is therefore unavailable at the depths of intimacy, and what is offered there is given on trust, into a history that will not run again, before its value can be confirmed.
\end{claim}

This structure is not a defect of intimate life to be engineered away, and naming it as a wager is not a counsel of resignation. It is the form that the generativity of relational value takes when the relation is a whole shared life, and it is inseparable from what makes intimacy the deepest of the domains in which worlds are crossed. The task of the following sections is to ask what the cultivation of these capacities can mean, once it is granted that at their depths they cannot be rehearsed, and to distinguish the capacities that can still be practiced in the ordinary way from those that can only be borne as a standing disposition into the wager.
